% ---
% Abstract
% ---
\begin{resumo}[Abstract]
 \begin{otherlanguage*}{english}
   The power factor is a fundamental indicator of energy efficiency in electrical systems, and its continuous monitoring is essential to avoid regulatory penalties and optimize energy consumption. This work proposes the development of an intelligent real-time electrical monitoring system capable of classifying different types of loads and estimating the power factor through the integration of low-cost sensing, Internet of Things (IoT)-based communication, and machine learning models. The system was implemented using the PZEM-004T sensor for acquiring electrical quantities, the NodeMCU ESP8266 microcontroller for data transmission via the MQTT protocol, and the Node-RED platform for real-time processing and visualization. The classification and regression models were developed based on the Multilayer Perceptron (MLP) architecture, trained on the Edge Impulse platform, and exported for execution in a WebAssembly environment. For experimental validation, 1000 samples were collected in a controlled laboratory environment for each of the eight considered load classes, totaling 8000 samples. The obtained results demonstrate the technical feasibility of the proposed solution for real-time load classification and power factor estimation. The developed system stands out for its low cost, scalability, and potential for application in real-world scenarios, contributing to the advancement of intelligent monitoring solutions in the context of energy efficiency.

   \textbf{Keywords}: Electrical monitoring; Power factor; Load classification; Artificial neural networks; Internet of Things.
 \end{otherlanguage*}
\end{resumo}